\section{Problem Formulation}

% General system model.
% Requires \usepackage{tikz} and \usetikzlibrary{arrows.meta, shapes.geometric} in the preamble.
%
% The mappings in the boxes (e.g. "x \mapsto D_0") are shown when \ifshowmappings is true.
% It is true by default. To hide them, put the following in the preamble:
%   \newif\ifshowmappings
%   \showmappingsfalse
\makeatletter
\@ifundefined{ifshowmappings}{\newif\ifshowmappings\showmappingstrue}{}
\makeatother
\providecommand{\sysmap}[1]{\ifshowmappings\\ $#1$\fi}
\begin{figure}[htbp]
\centering
\resizebox{\linewidth}{!}{%
\begin{tikzpicture}[
  box/.style={draw, rounded corners=2pt, align=center, minimum height=0.9cm, inner sep=3pt, font=\scriptsize},
  loop/.style={box, fill=blue!10, draw=blue!50!black},
  store/.style={draw, cylinder, shape border rotate=90, aspect=0.22, align=center, minimum height=1.2cm, inner sep=2pt, font=\scriptsize},
  sink/.style={draw=blue!50!black, ellipse, fill=blue!10, minimum height=0.6cm, inner sep=2pt, font=\scriptsize},
  arr/.style={-{Stealth[length=1.8mm]}, semithick},
  looparr/.style={arr, draw=blue!50!black},
  lbl/.style={font=\scriptsize},
  looplbl/.style={font=\scriptsize, text=blue!50!black}
]
  % Rounds t >= 1
  \draw[blue!50!black, dashed, fill=blue!3, rounded corners=3pt] (6.5,5.6) rectangle (15.5,7.25);
  \node[looplbl, anchor=north west] at (6.5,7.25) {Rounds $t \geq 1$};

  % Data stores
  \node[store, minimum width=1.4cm] (K) at (1.2,3.6) {Corpus\\ $\mathcal{K}$};
  \node[store, minimum width=2.1cm] (D) at (4.9,3.6) {Evidence store\\ $D_{t} \subseteq \mathcal{K}$};
  \node[store, minimum width=1.6cm] (lab) at (14.7,3.6) {Labelled\\ queries};

  % Round 0
  \node[box, minimum width=2.6cm] (ret) at (1.45,1.2) {Initial retrieval\\ over hops $h = 0, 1, \ldots$\sysmap{x \mapsto D_{0}}};
  \node[box, minimum width=2.2cm] (llm) at (4.5,1.2) {LLM\sysmap{(x, D_{0}) \mapsto y}};
  \node[box, minimum width=2.4cm] (ext) at (7.25,1.2) {Subclaim extraction\sysmap{y \mapsto \hat{y}_{0}}};
  \node[box, minimum width=2.6cm] (score) at (8.5,3.6) {Subclaim scoring\sysmap{(x, \hat{y}_{t}, D_{t}) \mapsto s_t}};
  \node[box, minimum width=2.3cm] (dec) at (12.0,3.6) {Decision\sysmap{s_t \mapsto (K_{t}, W_{t}, Z_{t})}};
  \node[box, minimum width=2.8cm] (sel) at (11.9,1.2) {Final selection\sysmap{(\hat{y}_{t^*}, D_{t^*}, s_{t^*}) \mapsto \hat{y}_{\mathrm{final}}}};

  % Rounds t >= 1 only
  \node[loop, minimum width=3.0cm] (rev) at (8.6,6.4) {Subclaim revision\sysmap{(x, K_{t}, W_{t}, D_{t+1}) \mapsto \hat{y}_{t+1}}};
  \node[loop, minimum width=2.4cm] (qgen) at (12.0,6.4) {Query generation\sysmap{(x, c, D_{t}) \mapsto q_t(c)}};
  \node[sink] (rm) at (14.4,6.4) {Removed};

  % Round 0 arrows
  \draw[arr] (-0.5,1.2) -- node[lbl, above] {$x$} (ret.west);
  \draw[arr] (-0.25,1.2) -- (-0.25,0.15) -- (llm.south |- 0,0.15) -- (llm.south);
  \draw[arr] ([xshift=-0.3cm]K.bottom |- ret.north) -- node[lbl, left] {$q_{h}$} ([xshift=-0.3cm]K.bottom);
  \draw[arr] ([xshift=0.3cm]K.bottom) -- node[lbl, right] {$\hat{D}_{h}$} ([xshift=0.3cm]K.bottom |- ret.north);
  \draw[arr] (ret.east) -- node[lbl, below] {$D_{0}$} (llm.west);
  \draw[arr] (3.0,1.2) -- (3.0,2.4) -- (D.bottom |- 0,2.4) -- (D.bottom);
  \draw[arr] (llm.east) -- node[lbl, above] {$y$} (ext.west);
  \draw[arr] ([xshift=-0.4cm]ext.north) |- node[lbl, left, pos=0.65] {$\hat{y}_{0}$} (score.west |- 0,3.35);
  \draw[arr] (D.east |- 0,3.85) -- node[lbl, above] {$D_{t}$} (score.west |- 0,3.85);
  \draw[arr] (score.east) -- node[lbl, above] {$s_t$} (dec.west);
  \draw[arr] (dec.south) -- node[lbl, left, align=right] {stop at time $t^*$\\ $\hat{y}_{t^*}, D_{t^*}, s_{t^*}$} (dec.south |- sel.north);
  \draw[arr] (sel.east) -- node[lbl, above] {$\hat{y}_{\mathrm{final}}$} (15.7,1.2);
  \draw[arr, dashed] (lab.bottom) |- (12.8,2.5) -- (12.8,0 |- sel.north);

  % Rounds t >= 1 arrows
  \draw[looparr] (dec.north) -- node[looplbl, left, align=center, pos=0.64] {rework\\ $W_{t}$} (qgen.south);
  \draw[looparr] ([xshift=-0.9cm]dec.north) |- node[looplbl, below, align=center, pos=0.85] {keep $K_{t}$,\\ rework $W_{t}$} (9.9,4.95) -- (9.9,0 |- rev.south);
  \draw[looparr] (8.55,0 |- rev.south) -- node[looplbl, left] {$\hat{y}_{t+1}$} (8.55,0 |- score.north);
  \draw[looparr] ([xshift=0.9cm]dec.north) |- node[looplbl, below, align=center, pos=0.75] {remove $Z_{t}$} (14.4,4.95) -- (14.4,0 |- rm.south);
  \draw[looparr] (qgen.north) -- (qgen.north |- 0,7.7) -- node[looplbl, above] {$q_t(c),\ c \in W_{t}$} (K.top |- 0,7.7) -- (K.top);
  \draw[looparr] (K.east) -- node[looplbl, above, yshift=1pt] {$\mathrm{Retrieve}(q_t(c))$} (D.west);
  \draw[looparr] (D.top) -- (D.top |- rev.west) -- node[looplbl, above, pos=0.35] {$D_{t+1}$} (rev.west);
  \draw[looparr, dashed] (lab.west) -- (dec.east);
\end{tikzpicture}%
}
\caption{Schema of the system. Round $0$, in black, is used by every method, and the rounds $t \geq 1$, in blue, are only used when the decision sends subclaims to be reworked. The dashed arrows show where the labelled queries are used for calibration.}
\label{fig:system}
\end{figure}

In all of the methods given in Section~\ref{sec:background}, the evidence used to guage factuality of the subclaims is retrieved in a single phase before the subclaims are generated. The subclaims themselves are also then only filtered once. The main consequence of this setup is that a true subclaim whose evidence was not retrieved has no way of being recovered and will be removed from the final set. In this work, we seek to alleviate this by allowing the system to consider its current subclaims, and either stop and answer, or retrieve further evidence for the subclaims that it cannot yet support and revise them. 

% We start by describing such a system by the schema in Figure~\ref{fig:system}, and explain how the methods above can be recovered with specific choices under this schema. 

The system answers a query $x$ over rounds $t = 0, 1, \ldots, T$, for a budget $T$ that is fixed in advance, and at round $t$ it holds the subclaims $\hat{y}_t$ and the evidence $D_t$. Round $0$, shown in black, is used by every method described in Section~\ref{sec:background}, whereas the rounds $t \geq 1$, shown in blue, are only reached by our system that reworks its subclaims.

Round $0$ starts with the initial retrieval, which maps the query $x$ to the evidence $D_0$. The system has access to a retriever over a fixed corpus $\mathcal{K}$, and we write $\mathrm{Retrieve}(q)$ for the set of documents that it returns for a retrieval query $q$. As in the case of multi-hop RAG, the initial retrieval may run over several hops $h = 0, 1, \ldots$, where at hop $h$ it sends a query $q_h$ to the retriever and receives the documents $\hat{D}_h = \mathrm{Retrieve}(q_h)$. We denote with $D_0$ the union of the documents retrieved over all hops. \tk{explain how the queries $q_h$ are formulated or not important?}

An LLM $G$ then writes the claim $y = G(x, D_0)$ from the query and the initial retrieved evidence. Subclaim extraction is performed by another LLM $\mathrm{Ext}$ \tk{label this and explain or not?} that splits $y$ into the subclaims $\hat{y}_0 = \mathrm{Ext}(y)$, as in the methods defined in Section~\ref{sec:background}. The evidence $D_0$ is also placed into an evidence store $\mathcal{E}$, which holds the evidence $D_t \subseteq \mathcal{K}$ collected up to round $t$.

At each round $t$, a subclaim scorer $\mathrm{Score}$ assigns a score $s_t(c)$ to every subclaim $c \in \hat{y}_t$ based on the query and the current retrieved evidence $D_t$, and we write $s_t = \mathrm{Score}(x, \hat{y}_t, D_t)$ for the scores of all subclaims of $\hat{y}_t$.

The decision rule $\mathrm{Dec}$ then uses the scores to partition $\hat{y}_t$ into three sets $(K_t, W_t, Z_t) = \mathrm{Dec}(s_t)$, where $K_t$ contains the subclaims that are retained, $W_t$ the subclaims that are reworked, and $Z_t$ the subclaims that are removed. It also decides whether to stop at the current round $t$, and we write $t^* \leq T$ for the round at which it does, and which may depend on the data. When it stops, the system passes the remaining subclaims $\hat{y}_{t^*} = K_{t^*} \cup W_{t^*}$, together with $D_{t^*}$ and $s_{t^*}$, to the final selection. A set of labelled queries, for which the labels $a(c)$ are known, is also available to calibrate both the decision and the final selection before the system is used to answer new queries, as shown by the dashed arrows in Figure~\ref{fig:system}.

Whilst the system has not stopped, it reworks the subclaims of $W_t$ in the rounds shown in blue. For each subclaim $c \in W_t$, query generation $\mathrm{Query}$ writes a retrieval query $q_t(c) = \mathrm{Query}(x, c, D_t)$ from the query, the subclaim and the evidence, and the documents returned for these queries are added to the evidence store,
\begin{equation}
  D_{t+1} = D_t \cup \bigcup_{c \in W_t} \mathrm{Retrieve}\bigl(q_t(c)\bigr).
  \label{eq:store}
\end{equation}
Subclaim revision $\mathrm{Rev}$ then gives the subclaims of the next round, $\hat{y}_{t+1} = \mathrm{Rev}(x, K_t, W_t, D_{t+1})$, in which the subclaims of $K_t$ remain the same, each subclaim of $W_t$ is rewritten given the new evidence or left unchanged. New subclaims may also be added, while the subclaims of $Z_t$ are removed. The next round  begins from the subclaim scoring of the updated subclaims $\hat{y}_{t+1}$.
% Each of these choices, as well as the decision to stop, may depend on $x$ and on everything observed in earlier rounds.

Finally, the final selection filter $\mathrm{Sel}$ returns the answer $\hat{y}_{\mathrm{final}} = \mathrm{Sel}(\hat{y}_{t^*}, D_{t^*}, s_{t^*})$ from the remaining subclaims, with $\hat{y}_{\mathrm{final}} \subseteq \hat{y}_{t^*}$.

In all three of the prior methods defined in Section~\ref{sec:background}, the decision keeps every subclaim, i.e., $K_0 = \hat{y}_0$ and $W_0 = Z_0 = \emptyset$, and the system therefore stops at $t^* = 0$. The final selection is then the filter $F_{\hat{\tau}}$ given in \eqref{eq:filter_set} applied to the subclaims of round $0$, i.e., $\hat{y}_{\mathrm{final}} = F_{\hat{\tau}}(\hat{y}_0)$, with $\hat{\tau}$ calibrated on the labelled queries. The three methods differ only in the initial retrieval and in the subclaim scorer. Conformal factuality uses no retrieval, thus $D_0 = \emptyset$, the LLM writes $y$ from $x$ alone, and $s_0(c)$ is the score $s(c)$ obtained from the LLM. Conformal-RAG uses a single hop with $q_0 = x$, and the method of Rehman and Yeh uses several hops, in which each query $q_h$ is rewritten by an LLM from $x$ and the documents received in earlier hops \tk{worth being more precise here?}. Both use the relevance score $R(c)$ of \eqref{eq:crag_score} on $D_0$ as $s_0(c)$. 

There also exists adaptive retrieval methods that retrieve during or after generation, which corresponds to the rounds $t \geq 1$, although they revise sentences of the answer rather than subclaims directly. For example, FLARE \cite{jiang2023active} retrieves whenever the next sentence of the answer contains low-confidence tokens and regenerates that sentence, and RARR \cite{gao2023rarr} retrieves evidence for a finished answer and revises the parts that the evidence does not support. Reworking subclaims in this way has therefore been considered before, but neither method gives a guarantee of the form \eqref{eq:cf}. \tk{This ought to go in the background but then it wouldn't be fair to refer to the diagram in a general sense as an augmentation of \emph{all} the methods in the background section. What do you thinl?}

The answer that the system finally returns must still satisfy the guarantee \eqref{eq:cf}, that is,
\begin{equation}
  \Pr\bigl[\, y^* \Rightarrow c \ \text{for every} \ c \in \hat{y}_{\mathrm{final}} \,\bigr] \geq 1 - \alpha,
  \label{eq:cf_seq}
\end{equation}
where the probability is now also taken over the choices that the system makes in each round. Since abstaining always satisfies the guarantee, we measure how useful the system is by the expected number of supported subclaims in its answer,
\begin{equation}
  U = \mathbb{E}\Bigl[ \sum_{c \in \hat{y}_{\mathrm{final}}} a(c) \Bigr],
  \label{eq:usefulness}
\end{equation}
and the goal is to ensure that $U$ as large as possible while keeping the guarantee \eqref{eq:cf_seq}.